\documentclass[12pt]{article}

%% Basic document formatting
\usepackage{amsmath, amsthm, amssymb, amsfonts}
\usepackage{mathrsfs}
\usepackage{mathtools}
\usepackage{xspace}
\usepackage{thmtools}
\usepackage{graphicx}
\usepackage{tikz}
\usepackage{tikz-cd} 
\usepackage{ytableau}
\usepackage{blkarray}
\usepackage{hyperref}
\usepackage{caption}
\usepackage{setspace}
\usepackage{array}
\usepackage{fancyhdr}
\usepackage{titling}
\usepackage[left=0.4in,right=0.4in,top=1in,bottom=1in]{geometry}
\usepackage{float}
\usepackage{cancel}
\usepackage{ulem}
\usepackage{tabularx}
\usepackage[utf8]{inputenc}
\usepackage[english]{babel}
\usepackage{framed}
\usepackage[dvipsnames]{xcolor}
\usepackage{environ}
\usepackage{tcolorbox}
\tcbuselibrary{theorems,skins,breakable}

\usepackage{enumitem}
\setlist[enumerate]{leftmargin=*}
\setlist[enumerate,1]{labelindent=\parindent}
\setlist[enumerate,2]{labelindent=0pt}

% Blackboard Bold
\newcommand{\Z}{\mathbb{Z}}                 % integers
\newcommand{\Zpl}{\mathbb{Z}_{+}}           % positive integers
\newcommand{\Znn}{\mathbb{Z}_{\geq 0}}      % nonnegative integers. 
\newcommand{\Q}{\mathbb{Q}}                 % rationals
\newcommand{\Qpl}{\mathbb{Q}_{+}}           % positive Rationals
\newcommand{\R}{\mathbb{R}}                 % reals
\newcommand{\Rpl}{\mathbb{R}_{+}}           % positive Reals
\newcommand{\C}{\mathbb{C}}                 % complex numbers
\newcommand{\F}{\mathbb{F}}                 % field

% Words
\newcommand{\st}{\text{ such that }}
\newcommand{\wrt}{\text{ with respect to }}
\newcommand{\with}{\text{ with }}
\newcommand{\ie}{\text{, i.e. }}
\newcommand*{\ditto}{\text{---}\texttt{"}\text{---}}

% Operators
\newcommand{\abs}[1]{\left|#1\right|}                   % absolute value
\newcommand{\norm}[1]{\abs{\abs{#1}}}
\newcommand{\floor}[1]{\left\lfloor #1 \right\rfloor}   % floor
\newcommand{\ceil}[1]{\left\lceil #1 \right\rceil}      % ceiling
\newcommand{\powset}[1]{\mathscr{P}(#1)}

% Category Theory 
\renewcommand{\hom}{\text{Hom}}
\newcommand{\homspace}[3]{\hom_{#1}(#2, #3)}
\newcommand{\coproduct}[9]{
  \begin{tikzcd}[ampersand replacement=\&]
      \& \& #1 \arrow[dll, bend right, "#6"] \arrow[dl, "#5"] \\
   #4 \& #3 \arrow[l, "#9"] \\
      \& \& #2 \arrow[ull, bend left, "#8"] \arrow[ul, "#7"]
  \end{tikzcd}
}

\newcommand{\product}[9]{ 
  \begin{tikzcd}[ampersand replacement=\&]
      \& \& #3 \\
      #1 \arrow[r, "#7"] \arrow[urr, bend left, "#5"] \arrow[drr, bend right, "#6"] \& #2 \arrow[ur, "#8"] \arrow[dr, "#9"] \\ 
      \& \& #4
  \end{tikzcd}
}


% Algebra 
\newcommand{\Syl}[2]{\text{Syl}_{#1}{(#2)}}           % set of Sylow-p subgroups 
\newcommand{\<}{\langle}                % \<x\>, subgroup generated by x 
\renewcommand{\>}{\rangle}
\renewcommand{\index}[2]{[#1 : #2]}
\newcommand{\id}{\text{id}}             % identity element
\newcommand{\lhdneq}{%
  \mathrel{\ooalign{$\lneq$\cr\raise.22ex\hbox{$\lhd$}\cr}}}
\newcommand{\order}[1]{\text{o}(#1)}    % order of an element
\newcommand{\spec}{\operatorname{Spec}}
\newcommand{\rad}{\operatorname{rad}}   % radical of an ideal 
\newcommand{\sym}[1]{\mathfrak{S}_{#1}}
\newcommand{\aut}[1]{\text{Aut}(#1)} 
\newcommand{\gal}[2]{\text{Gal}(#1 / #2)} 
\newcommand{\inn}[1]{\text{Inn}(#1)} 
\let\oldcong\cong
\let\oldequiv\equiv
\renewcommand{\cong}{\oldequiv}
\renewcommand{\equiv}{\oldcong}

\newcommand{\triangleleftneq}{\mathrel{\ooalign{%
  $\trianglelefteq$\cr
  \hidewidth\raise0.06ex\hbox{$\not\mkern4mu$}\hidewidth\cr
}}}

\makeatletter % cycle
\newcommand{\cyc}[1]{(\mathbf{\cyc@process#1\relax})}
\def\cyc@process#1#2\relax{%
	#1%
	\ifx\relax#2\relax
	\else
		\,\cyc@process#2\relax
	\fi
}
\makeatother

\newcommand{\GL}[2]{\text{GL}_{#1}(#2)}                             % general linear group
\newcommand{\SL}[2]{\text{SL}_{#1}(#2)}                             % special linear group
\newcommand{\gl}[2]{\mathfrak{gl}_{#1}(#2)}
\renewcommand{\sl}[2]{\mathfrak{sl}_{#1}(#2)}
\newcommand{\so}[2]{\mathfrak{so}_{#1}(#2)}
\newcommand{\su}[1]{\mathfrak{su}(#1)}

% Linear Algebra
\newcommand{\bmat}[1]{\begin{bmatrix}#1\end{bmatrix}}   % bracketed matrix
\newcommand{\rank}{\operatorname{rank}}                 % rank
\newcommand{\nullity}{\operatorname{nullity}}           % nullity
\newcommand{\tr}{\operatorname{tr}}

% Combinatorics
\newcommand{\qnum}[1]{\left[#1\right]_q}                % q-analog
\newcommand{\qfact}[1]{\left[#1\right]!_q}              % q-analog factorial 
\newcommand{\qbinom}[2]{\binom{#1}{#2}_q}               % Gaussian binomial coefficient

% Topology/Analysis
\newcommand{\ball}[2]{\text{B}_{#1}(#2)}  % B_r(x): open r-balls around x
\newcommand{\diam}{\text{diam}}           % diamter of a set in metric space 
\newcommand{\lowint}[2]{\underline{\int_{#1}^{#2}}}
\newcommand{\upint}[2]{\overline{\int}_{#1}^{#2}}

% Misc Notation
\newcommand{\oneton}[1]{\{1, \ldots, #1\}}
\newcommand{\defeq}{\vcentcolon=}   % :=
\newcommand{\eqdef}{=\vcentcolon}   % =: 
\renewcommand{\bf}[1]{\textbf{#1}}
\newcommand{\im}{\operatorname{im}}
\newcommand{\fnrest}[2]{\left.#1\right|_{#2}}
\newcommand{\isom}{\overset{\sim}{\rightarrow}} 


\renewcommand{\thesection}{\S \arabic{section}}
\renewcommand{\thesubsection}{\arabic{subsection}.}
\renewcommand{\thesubsubsection}{(\alph{subsubsection})} 

\newtheorem{claim}{Claim}
\newtheorem*{lemma}{Lemma}

%% Headers & title setup
\newcommand{\course}{\textit{Algebra: Chapter 0}, Aluffi} 
\newcommand{\myname}{Jay Ser}
\setlength{\headheight}{14.5pt}
\pagestyle{fancy}
\fancyhf{}
\renewcommand{\headrulewidth}{0.4pt}
\lhead{\course}
\rhead{\myname}
\cfoot{\thepage}
\setlength{\droptitle}{-4em} 
\title{\course\ \\ Part 5: Irreducibility and Factorization in Integral Domains}
\author{\myname}
\date{May 6, 2026 $\sim$ \today}

\begin{document}
\maketitle
\thispagestyle{fancy}

\newcommand{\dirsum}[2]{#1^{\oplus#2}}
\newcommand{\rmod}[1]{#1\text{-module}\xspace}
\newcommand{\idealp}{\mathfrak{p}}
\newcommand{\idealq}{\mathfrak{q}}
\newcommand{\idealm}{\mathfrak{m}} 

%------------------------------------------------------------------------------%

\section{Chain Conditions and Existence of Factorizations} 
\setcounter{subsection}{2} 
\subsection{} 
\subsubsection{} 
The action of $k[t]$ on $k[x]$ is determined by the ring homomorphisms $k[t] \overset{\varphi}{\rightarrow} R \hookrightarrow k[x]$; 
consequently, the action is entirely by the action of constants in $k$ and the indeterminant $t$: 
\begin{align*}
  a \cdot g(x) & = ag(x) \hspace{2em} (a \in k) \\ 
  t \cdot g(x) & = f(x) g(x) 
\end{align*}
Let $f(x) = a_n x^n + \ldots + a_0$. 
The definitions above indicate that $1, x, \ldots, x^{n - 1}$ generate $k[x]$ as a $\rmod{k[t]}$: 
as a ring, $k[x]$ is generated by $1, x, x^2, \ldots$, so it suffices to show that one can retrieve all powers of $x$ by the action of $k[t]$ on polynomials in $x$ of degree less than $n$. 
Indeed, 
\begin{equation} \label{eq:1_3_xpowers}
  x^n = a_{n}^{-1} \cdot (t \cdot 1 - a_{n - 1} \cdot x^{n - 1} - \ldots - a_0 \cdot 1) 
\end{equation}
shows $x^n \in \<1, x, \ldots, x^{n - 1}\>$, and induction and a similar formula as \eqref{eq:1_3_xpowers} shows all powers of $x$ are generated by $\<1, x, \ldots, x^{n - 1}\>$. 

\subsubsection{}
$k[x]$ is a finitely generated $\rmod{k[t]}$ and $k[t]$ is a Noetherian ring (not yet covered in the book;
polynomial rings over fields are PIDs), 
so $k[x]$ is a Noetherian $\rmod{k[t]}$. 
Hence to prove $R$ is finitely generated, it suffices to check that $R$ is a $k[t]$-submodule of $k[x]$. 
Indeed, because $R$ contains the element $f$ and field $k$, for any $p(t) \in k[t]$ and any $g[x] \in R$, 
$$p(t) \cdot g = p(f) g \in R.$$

\subsubsection{} 
$k$ itself is a field, so it is a Noetherian ring. 
Take a subring of $k[x]$ that properly contains $k$; 
that is, consider $R \subset k[x]$ as in the last two parts, its $\rmod{k[t]}$ structure induced by some nonconstant polynomial $f \in R$. 
By part (b), $R$ is a Noetherian $\rmod{k[t]}$. 
Take an ideal $I \subset R$. 
It is clear that $I$ is a $\rmod{k[t]}$ of $R$, so $I$ is finitely generated as a $\rmod{k[t]}$. 
Suppose $I = \<a_1, \ldots, a_n\>$ as a $\rmod{k[t]}$. 
For all $a \in I$, there exists $p_1(t), \ldots, p_n(t) \in k[t]$ such that 
\begin{align*} 
  a & = p_1(t) \cdot a_1 + \ldots + p_n(t) \cdot a_n \\ 
    & = p_1(f)a_1 + \ldots + p_n(f)a_n
\end{align*}
Notice that the third equality expresses $a$ as a $R$-linear combination of $a_1, \ldots, a_n$ (recall the action of $k[t]$ on $R$; $p_1(f), \ldots, p_n(f)$ are elements of $R$). 
Consequently, $I = \<a_1, \ldots, a_n\>$ as an $R$-module; 
that is, $I = \<a_1, \ldots, a_n\>$ as an ideal in $R$, hence $I$ is finitely generated. 
This proves $R$ is Noetherian. 

\setcounter{subsection}{7}
\subsection{} 
Let $\mathscr{F}$ be as defined in the hint, and assuming $\mathscr{F} \neq \emptyset$, let $I_0$ be the maximal element of $\mathscr{F}$. 
Since $I_0$ is not prime, find $a, b \notin I_0$ such that $ab \in I_0$. 
Since $A \defeq \<I_0, a\>$ and $B \defeq \<I_0, b\>$ properly contain $I_0$, $A$ and $B$ both contain a product of prime ideals; 
say $P_1 \cdots P_t \subset A$ and $Q_1 \cdots Q_s \subset B$ for prime ideals $P_i$ and $Q_j$. 
Then notice 
$$AB = \<I_0, ab\> = \<I_0\>$$ 
and 
$$P_1 \cdots P_t Q_1 \cdots Q_s \subset AB,$$
hence $I_0$ contains a product of prime ideals, contradicting the construction of $\mathscr{F}$. 
It follows $\mathscr{F} = \emptyset$.

\setcounter{subsection}{9}
\subsection{} 
It is clear that $R / I$ is Artinian due to the inclusion-preserving bijection between ideals in $R$ containing $I$ and ideals in $R / I$. 
If $R$ is an Artinian integral domain, the chain 
$$\<a\> \supset \<a^2\> \supset \<a^3\> \supset \ldots$$ 
for any $a \in R \setminus 0$ shows that $\<a^n\> = \<a^{n + 1}\>$ for some $n$, 
i.e., $a^n = u a^{n + 1}$ for a unit $u \in R$. 
The cancellation law applies because $R$ is an integral domain, giving $1 = ua$.
$u$ is the inverse of $a$, thus proving that $R$ is a field. 

To show that an Artinian ring has Krull dimension 0, one must show if $\mathfrak{p}, \mathfrak{p}'$ are prime ideals in $R$, 
$$\mathfrak{p} \subset \mathfrak{p}' \implies \mathfrak{p} = \mathfrak{p}'.$$ 
Indeed, one sees this is so by mapping $\idealp$ and $\idealp'$ to $R / \idealp$.
Using bar notation, $\overline{\idealp} = \overline{0}$ and $\overline{\idealp'}$ is a prime ideal in the quotient $R / \idealp$. 
But the quotient ring is a field: that $\idealp$ is prime means $R / \idealp$ is an integral domain, and $R / \idealp$ Artinian. 
Fields have no nontrivial, nonproper ideals, and since $\overline{\idealp} \neq R / \idealp$, it follows $\overline{\idealp'} = \overline{0}$. 
This proves $\idealp' = \idealp$, as desired. 

\section{UFDs, PIDs, Euclidean Domains} 
\setcounter{subsection}{5}
\subsection{} 
For any $a, b \in R \setminus 0$, there exists $c \in R \st \<c\> = \<a\> \cap \<b\>$. 
By definition of the least common multiple, $c$ is the LCM of $a$ and $b$: 
if for any $m \in R$, $\<m\> \subset \<a\>$ and $\<m\> \subset \<b\>$, then $\<m\> \subset \<c\>$. 

$ab / c$ is then the GCD of $a$ and $b$.
Suppose $\<a\> \subset \<e\>$ and $\<b\> \subset \<e\>$. 
$ab / e$ is then a common multiple of $a$ and $b$, so $\<ab / e\> \subset \<c\>$. 
In other words, for some $q \in R$, 
$$cq = ab / e \implies ceq = ab \implies eq = ab / c \implies e \mid ab/c,$$
verifying $ab / c$ is the gcd of $a$ and $b$. 

\subsection{} 
For any elements $a, b \in R \setminus 0$, let $d = ra + sb$ be the GCD of $a$ and $b$. 
Namely, the fact that $d$ is a $R$-linear combination means  
\begin{equation} \label{eq:2_7_pd}
  \<a, b\> = \<d\>.
\end{equation}
$R$ being Noetherian means every ideal $I \lhd R$ is finitely generated. 
But by \eqref{eq:2_7_pd}, one can induct on the number of generators to show that any finitely generated ideal is in in fact generated by a single element;
this shows $R$ is a PID. 

\setcounter{subsection}{8} 
\subsection{}
Let $\idealp$ be a prime ideal in a UFD $R$ of height 1. 
Namely, since 0 is a prime ideal in any integral domain, the height of $\idealp$ says $\idealp$ does not properly contain any nonzero prime ideal of $R$. 
Take any $a \in \idealp \setminus 0$. 
Let $a = p_1\cdots p_k$ be the prime decomposition for $a$. 
Since $\idealp$ is prime, $a \in \idealp$ means some $p_i \in \idealp$, i.e., 
$$\<p_i\> \subset \idealp.$$
By the height of $\idealp$, this inclusion must, in fact, be an equality. 
Hence $\idealp$ is principal.

\setcounter{subsection}{11} 
\subsection{} 
Let $R[x]$ be a PID and take any $a \in R \setminus 0$. 
By assumption, there exists $f(x) \in R[x]$ such that 
\begin{equation} \label{eq:2_12_pric}
  \<a, x\> = \<f\>.
\end{equation}
\eqref{eq:2_12_pric} says $f \mid x$, but because $R[x]$ is an integral domain, necessarily $f = 1$ or $f = x$ (up to units). 
Furthermore, $f \mid a$ means $f \neq x$ since, once again, $R[x]$ is an integral domain. 
So $f = 1$, meaning $\<a, x\> = \<1\>$. 
Accordingly, there are $g(x), h(x) \in R[x]$ such that 
$$ag + xh = 1$$
Namely, $xh$ is a polynomial with no constant term, so, letting $c$ be the constant term of $g$, 
$$ac = 1.$$
Namely, this shows that $a$ is a unit in $R$. 
This proves $R$ is a field. 

\subsection{} 
The proof of $x^a - 1\mid x^b - 1 \iff a \mid b$ follows the same way as the proof of $c^a - 1 \mid c^b - 1 \iff a \mid b$; 
I will only show the latter. 
If $a \mid b$, then $b = aq$ for some $q$. 
$$c^{b} - 1 = c^{aq} - 1 = (c^a - 1)(c^{a(q - 1)} + c^{a(q - 2)} + \ldots + c^a + 1)$$
shows $c^a - 1 \mid c^b - 1$. 

Conversely, suppose $c^a - 1 \mid c^b - 1$. 
Write $b = aq + r$ for some $0 \leq r < a$. 
$c^a - 1 \mid c^{aq} - 1$ always, so 
$$c^a - 1 \mid [c^{aq}c^r - 1 - (c^{aq} - 1)] = c^{aq}(c^r - 1)$$
Since $c^a - 1 \mid c^{aq - 1}$, $\gcd(c^a - 1, c^{aq}) = 1$, which forces 
\begin{equation} \label{eq:2_13_div}
  c^a - 1 \mid c^r - 1. 
\end{equation}
But since $r < a$, \eqref{eq:2_13_div} is nonsense unless $c^r - 1 = 0$, i.e., unless $r = 0$. 
This proves $a \mid b$. 

\setcounter{subsection}{16} 
\subsection{} 
Take any nonunit $c_1 \in R \setminus 0$. 
If $c_1$ is not the desired element, there is an element $a_1 \in R$ such that $a_1 = c_1 q_1 + r_1$ for some nonunit $r_1$ such that $r_1 \neq 0$, $v(r_1) < v(c_1)$. 
Let $c_2 = r_1$, and once again, if $c_2$ is not the desired element, find $a_2 \in R$ such that $a_2 = c_2 q_2 + r_2$ for some nonunit, nonzero $r_2$ with $v(r_2) < v(c_2)$, and repeat the process for $c_3 = r_2$, etc. 
This process must eventually terminate, yielding some $c_k$ such that for all $a \in R$, there is some $r, q \in R$ such that $a = c_k q + r$ and $r$ is zero or a unit. 
If not, since $v(r_i) < v(c_i)$ for each $c_i$ that doesn't satisfy the desired property, the $v$ will eventually start mapping to the negative numbers, which contradicts the fact that the codomain of $v$ is the nonnegative integers. 

\newcommand{\m}{\mathfrak{m}}
\section{Intermezzo: Zorn's Lemma} 
\setcounter{subsection}{10} 
\subsection{}  
Suppose every nonzero prime ideal is a maximal ideal. 
Then every maximal ideal is a principal ideal, which in turn shows that every prime ideal is a principal ideal: 
let $\m$ be a maximal ideal and pick any $a \in \m$. 
Let $a = q_1 \cdots q_k$ be the irreducible decomposition of $a$. 
Since $\m$ is a maximal ideal (namely, it is prime), some $q_i \in \m$. 
But irreducible elements are prime elements in a UFD and prime elements generate a maximal ideal by the hypothesis on $R$, hence $\<q_i\> = \m$. 
This shows that maximal ideal are principal. 

Let 
$$\mathcal{F} = \{I \neq 0 \lhdneq R \mid I \text{ is not a principal ideal}\}.$$
One can easily check that every chain in $\mathcal{F}$ has an upper bound in $\mathcal{F}$; 
by Zorn's lemma, $\mathcal{F}$ has a maximal element $I_0$. 
Since $I_0$ is not principal, it is necessarily not prime. 
Take $a, b \notin I_0$ such that $ab \in I_0$.
$I_0 \subsetneq \<I_0, a\>$ and the maximality of $I_0$ in $\mathcal{F}$ means $\<I_0, a\> = \<a'\>$ for some $a' \in R$. 
Similarly, $\<I_0, b\> = \<b'\>$ for some $b' \in R$. 
But then 
$$\<a' b'\> = \<a'\> \<b'\> = \<I_0, ab\> = \<I_0\>,$$
which contradicts the assumption that $I_0 \in \mathcal{F}$. 

\section{Unique Factorization in Polynomial Rings}
\setcounter{subsection}{3} 
\subsection{} 
The first statement is a tautology since every prime ideal in a PID is a prime principal ideal. 
However, primitive doesn't imply very primitive in a general UFD. 
Consider $R = \Z[x_1]$ and $f = 2x_1 x_2 + 5 \in R[x_2]$. 
Then $\gcd(2x_1, 5) = 1$ shows $f$ is primitive, but $2 \notin \<2x_1, 5\>$ shows $f$ is not very primitive. 

\setcounter{subsection}{7} 
\subsection{} 
Clearly $\sim$ is symmetric and reflexive. 
To check that it is transitive, suppose 
\begin{gather*} 
  t_1(s'm - sm') \\ 
  t_2(s''m' - s'm'') 
\end{gather*}
Then 
\begin{align*}
  t_1 t_2 s' (s''m - sm'') & = t_1 t_2 (s'' s' m - s'' sm' + s'' sm' - s' sm'') \\ 
						   & = 0
\end{align*}
shows $(m, s) \sim (m', s')$ and $(m', s') \sim (m'', s'')$ imply $(m, s) \sim (m'', s'')$. 

The natural action of $S^{-1}R$ on $S^{-1} M$ that makes the latter a $S^{-1}R$-module is 
$$r/s \cdot m/t \defeq (rm)/(st),$$
where $rm$ is, of course, induced by the $R$-module structure of $M$. 
It is easy to check that this operation satisfies the axioms of a $S^{-1}R$ module, should the operation be well defined. 
Indeed, it is: 
for example, suppose $r_1 / s_1 = r_2 / s_2$ in $S^{-1}R$ with $u(s_2 r_1 - s_1 r_2) = 0$ for some $u \in S$. 
It must be shown that $r_1 / s_1 \cdot m / t = r_2 / s_2 \cdot m / t$ for any $m / t \in S^{-1} M$. 
This is verified by the computation 
\begin{align*}
  u(r_1 s_2 t \cdot m - r_2 s_1 t \cdot m) & = tu(s_2 r_1 - s_1 r_2) \cdot m \\ 
										   & = 0 \cdot m = 0.
\end{align*}

\subsection{} 
\subsubsection{} 
Take $a_1 / s_1, a_2 / s_2 \in I^{e}$ and any $r / t \in S^{-1} R$. 
Then 
\begin{equation} \label{eq:4_9_a} 
  \frac{a_1}{s_1} + \frac{r}{t} \frac{a_2}{s_2} = \frac{ta_1 s_2 + ra_2 s_1}{ts_1 s_2}
\end{equation}
Since $I$ is an ideal, $ta_1 s_2 + r a_2 s_1 \in I$, so \eqref{eq:4_9_a} is in $I^{e}$. 
This shows $I^{e}$ is an ideal in $S^{-1}R$. 

To verify that $I^{e}$ is not proper, suppose $1 \in I^{e}$, i.e., 
there is exists some $a / s \in I^{e}$ and some $u \in S$ such that 
$$u(a - s) = 0.$$
In other words, $au = us$. 
But $au \in I$ while $us \in S$, contradicting the assumption $I \cap S = \emptyset$. 

\subsubsection{} 
Take $a, b \in J^{c}$ and $r \in R$. 
Since $J$ is an ideal, $(a + rb) / 1$ in $S^{-1}R$ is an element of $J$, hence $a + rb \in J^{c}$. 
Further, since $J$ is a proper ideal, $S \cap J^{c} = \emptyset$; for if there was some $s \in S \cap J^{c}$, then $J$ contains the unit $s$ in $S^{-1}R$, hence $J = \<1\>$. 

\subsubsection{} 
$(J^{c})^{e} = J$ is clear; 
see Q14 for an argument similar to proving $(J^{c})^{e} = J$. 

Take $a \in (I^{e})^{c}$. 
Then $a/1 \in I^{e}$, so $\exists a' / s \in I^{e}$ and $u \in S \st$
$$u(sa - a') = 0$$
which means $asu = a'u$, with the right hand side in $I$. 
Of course, $su \in S$, so $a$ satisfies the desired property. 

Vice versa, suppose $a \in R$ such that $sa \in I$ for some $s \in S$.
Then $as / s = a \in I^{e}$, hence $a \in (I^{e})^{c}$. 

This gives the relation $I \subset (I^{e})^{c}$. 

\subsubsection{} 
As given in the hint, let $R = \C[x, y]$ and $S = \{1, x, x^2, \ldots \} \subset R$. 
If $I = \<xy\>$, then by (c), $y \in (I^{e})^{c}$, but $y \notin I$ clearly. 
This phenemonon occurs because $I^{e}$ contains elements of the form 
$$xy f(x, y) / x^k, \quad (f(x, y) \in R, \, k \in \Znn),$$
and hence for $k = 1$, the $x$'s in the denominator and numerator cancel out. 

\subsection{} 
It is clear that if $\idealq \lhd S^{-1}R$ is prime, then its contraction $\idealq^{c} \lhd R$ is prime. 
Suppose $\idealp \lhd R$ is prime disjoint from $S$ and take $\frac{a}{s} \frac{b}{t} \in \idealp^{e}$ where $b / t \notin \idealp^{e}$. 
Then there exists $u / v \in \idealp^{e}$ with $u \in \idealp$ and $v \in S$ such that $u/v = (ab)/(st)$, i.e., there exists $k \in S$ such that 
$$kstu = kvab$$ 
in $R$.
The left hand side is an element of $\idealp$ since $u \in \idealp$. 
On the right hand side, $kv \in S$ means $kv \notin \idealp$, and further $b \notin \idealp$ since $b / t \notin \idealp^{e}$. 
Hence $kvb \notin \idealp$, which forces $a \in \idealp$ by primality of $\idealp$. 
This means $a /s \in \idealp^{e}$, proving the primality of $\idealp^{e}$. 

It is clear that the assignment $\idealp \mapsto \idealp^{e}$ preserves inclusions. 
The paragraph above verifies that this assignment indeed maps prime ideals of $R$ disjoint from $S$ to prime ideals of $S^{-1}R$, and its inverse map $\idealq \mapsto \idealq^{c}$ maps prime ideals of $S^{-1}R$ to prime ideals of $R$ disjoint from $S$. 
To see that this mapping is a bijection, notice that if $\idealp \lhd R$ is a prime ideal disjoint from $S$, then $(\idealp^{e})^{c} = \idealp$: 
as noted in Q9(c), $\idealp \subset (\idealp^{e})^{c}$. 
To go the other direction, take any $r \in (\idealp^{e})^{c}$, i.e., 
$\exists s \in S \st rs \in \idealp$. 
Since $\idealp$ is disjoint from $S$ and $\idealp$ is prime, $r \in \idealp$. 

\subsection{} 
Let $\idealp \lhd R$ be a prime ideal and define $S \defeq R \setminus \idealp$. 
Then $S$ is a multiplicative set: 
$1 \notin \idealp$, and $s_1, s_2 \notin \idealp \implies s_1 s_2 \notin \idealp$; 
in other words, $1 \in S$ and $s_1, s_2 \in S \implies s_1 s_2 \in S$. 

The inclusion-preserving bijection between prime ideals of $R$ contained in $\idealp$ and prime ideals of $R_{\idealp}$ is just the bijection established in Q10;
the ``prime ideals of $R$ contained in $\idealp$" detail is due to the definition $S = R \setminus \idealp$; 
hence any ideal contained in $\idealp$ is disjoint from $S$. 

Of course, $\idealp$ is the unique maximal element amongst prime ideals in $R$ contained in $\idealp$. 
Since the bijection established above is inclusion-preserving, $\idealp^{e} \lhd R_{\idealp}$ is the unique maximal ideal of $R$ (recall maximal ideals are prime ideals).  

\subsection{} 
Numbering the statements as 
\begin{enumerate} 
  \item $M = 0$ 
  \item $M_{\idealp} = 0$ for every prime ideal $\idealp \lhd R$. 
  \item $M_{\idealm} = 0$ for every maximal ideal $\idealm \lhd R$,
\end{enumerate} 
it is clear $(1) \implies (2) \implies (3)$ since maximal ideals are prime ideals. 
To prove the contrapositive of $(3) \implies (1)$, suppose $M$ is not the zero $R$-module and take any element $m \in M \setminus 0$. 
Let $I = \{a \in R \mid am = 0\}$. 
It is clear that $I$ is an ideal in $R$. 
Let $\idealm$ be a maximal ideal in $R$ containing $I$. 
Then $M_{\idealm}$ is not the zero module: 
on the contrary, if $M_{\idealm}$ is the zero module, then $m / s$ would be 0 for all $s \in R \setminus \idealm$. 
In other words, 
$$\exists t \in R \setminus \idealm \st t(m - 0) = 0,$$ 
which contradicts the assumption that $t \notin I$ and $I$ is precisely the ideal containing all the elements of $R$ that annihilate $m$. 

\setcounter{subsection}{13} 
\subsection{} 
Note that these extensions and contractions are with respect to $M \rightarrow S^{-1}M$, where $M$ is viewed as a $R$-module and $S^{-1}M$ is viewed as a $S^{-1}R$ module. 
Given a submodule $N \lhd M$, one can check 
$$N^{e} \defeq \{n / s \mid n \in N, \, s \in S\}$$ 
is a submodule of $S^{-1}M$. 
Similarly, given a submodule $\hat{N} \lhd S^{-1}M$, 
$$\hat{N}^{c} \defeq l^{-1}(\hat{N})$$
is a submodule of $M$ 
$l$ is the map $l: M \rightarrow S^{-1}M; \quad m \mapsto m/1$. 

Let's verify $\hat{N} = (\hat{N}^{c})^{e}$. 
Take any $n / s \in \hat{N}$. 
Then $\frac{s}{1} \cdot \frac{n}{s} = n / 1 \in \hat{N}$, so $n \in \hat{N}^{c}$. 
It immediately follows $n /s \in (\hat{N}^{c})^{e}$. 
Next, take any $n / s \in (\hat{N}^{c})^{e}$ where $n \in \hat{N}^{c}$ by definition. 
Then, one again by definition, $l(n) = n / 1 \in \hat{N}$, so $\frac{1}{s} \cdot \frac{n}{1} = n / s \in \hat{N}$. 

The result $\hat{N} = (\hat{N}^{c})^{e}$ indicates every localization of a Noetherian module is Noetherian. 
Assume $M$ is a Noetherian $R$-module and take a submodule $\hat{N} \lhd S^{-1}M$. 
$\hat{N}^{c}$ is a submodule of $M$, so it is finitely generated. 
This means all possible numerators of $(\hat{N}^{c})^{e} = \hat{N}$ are generated by finitely many elements; 
the action of $S^{-1}R$ on $S^{-1}M$ gives all the possible denominators of $\hat{N}$, so $\hat{N}$ is finitely generated. 

\subsection{} 
\subsubsection{} 
Let $q$ be irreducible in $R$ and let $q / 1 = \frac{r_1}{s_1} \frac{r_2}{s_2}$. 
Since $R$ is a UFD and particularly an integral domain, this equality in $S^{-1}R$ is equivalent to the equality 
$$s_1 s_2 q = r_1 r_2$$ 
in $R$. 
$q \mid r_1 r_2$ and irreducible elements are prime in UFDs, so $q \mid r_1$ or $q \mid r_2$. 
Without loss of generality, let $q \mid r_1$ and write $r_1 = q r_{1}'$. 
Hence 
\begin{align*}
  q / 1 & = \frac{qr_{1}' r_2}{s_1 s_2}  \\ 
  \implies 1 & = \frac{r_{1}' r_2}{s_1 s_2} 
\end{align*}
where the $q / 1$ can be canceled because $R$ is an integral domain $\implies$ $S^{-1}R$ is an integral domain. 
This shows that $r_2 / s_2$ is a unit, from which follows $q / 1$ is a unit or an irreducible element. 
It is indeed possible that $q / 1$ is a unit in $S^{-1}R$; 
for example, if $q \in S$. 

\subsubsection{}
Suppose $a / s$ is irreducible in $S^{-1}R$. 
Since $1 / s$ is a unit in $S^{-1}R$, $a / 1$ is irreducible in $S^{-1}R$, and one may as well show that $a / 1$ is associate to $q / 1$ for some $q \in R$ irreducible. 
If $a$ is irreducible in $R$, there is nothing more to be shown. 
If $a$ is not irreducible in $R$, let 
$$a = p_1 p_2 \cdots p_k$$
be the irreducible factorization of $a$. 
This relation holds in $S^{-1}R$ as well. 
However, since $a$ is irreducible in $S^{-1}R$, either $p_1$ or $p_2 \cdots p_k$ is a unit in $S^{-1}R$. 
If $p_2 \cdots p_k$ is a unit, then $a$ is associate to $p_1 / 1$ with $p_1$ irreducible in $R$. 
If $p_1$ is a unit, then $p_2 \cdots p_k$ is irreducible in $S^{-1}R$ and thus either $p_2$ or $p_3 \cdots p_k$ is a unit in $S^{-1}R$, and one can proceed as before to find that $a$ is associate to some $p_i / 1$. 

\subsubsection{} 
Take any ascending chain 
$$\<r_1 / s_1\> \subset \<r_2 / s_2\> \subset \cdots$$ 
in $S^{-1}R$. 
Since the $1 / s_i$'s are all units, the chain is above is equal to $\<r_1 / 1\> \subset \<r_2 / 1\> \subset \cdots$. 
In $R$, the chain $\<r_1\> \subset \<r_2\> \subset \cdots$ stabilizes, say at index $k$ (one must not forget that now the $\<r_i\>$'s are ideals in $R$). 
So consider $\<r_k / 1\>$ and $r_{k + j} / 1\>$ in $S^{-1}R$. 
Take any $\frac{a}{t}r_{k + j} \in \<r_{k + j} / 1\>$. 
Since the ideals $\<r_{k}\> = \<r_{k + j}\>$ in $R$, there is some $b \in R$ such that $br_{k} = ar_{k + j}$. 
Hence 
$$\frac{a}{t}r_{k + j} = \frac{b}{t} r_k \in \<r_k / 1\>,$$ 
i.e., $\<r_{k} / 1\> = \<r_{k + j}\>$. 
This verifies that the ascending chain condition for principal ideals is satisfied in $S^{-1}R$. 

Next, suppose that $a / s$ is irreducible in $S^{-1}R$ and that $a /s \mid \frac{r_1}{t_1} \frac{r_2}{t_2}$. 
By part (b), $a / s$ is associate to some $q / 1$ for $q$ irreducible in $R$ (of course, $q / 1$ is irreducible in $S^{-1}R$). 
To check that $a / s$ is prime, it suffices to check that $q / 1$ divides $r_1 / 1$ or $r_2 / 1$. 
But indeed, $q$ does divide $r_1$ or $r_2$ in $R$ because $R$ is a UFD and irreducibles are prime in a UFD; 
the result therefore holds in $S^{-1}R$ as well. 
This proves that $a / s$ is prime in $S^{-1}R$. 

Since the ascending chain condition and irreducible = prime condition hold in $S^{-1}R$, it is a UFD. 

\subsection{} 
Q15 already shows $S^{-1}R$ is a UFD if $R$ is a UFD. 
Vice versa, suppose $R$ is Noetherian and $S^{-1}R$ is a UFD. 
Recall that localizations of Noetherian rings are themselves Noetherian. 
Take a prime ideal $\idealp \lhd R$ of height 1. 
Suppose $\idealp \cap S \neq \emptyset$. 
Then $s^k \in \idealp$ for some power $k$. 
Since $\idealp$ is prime, $s \in \idealp$ and therefore $\idealp$ contains a nonzero prime ideal $\<s\>$ (since $s$ is a prime element). 
$\idealp$ is of height 1, so $\idealp = \<s\>$. 

Next, assume $\idealp \cap S = \emptyset$. 
$\idealp$ corresponds to the prime ideal $\idealp^{e} \lhd S^{-1}R$. 
By the inclusion-preserving bijection between prime ideals of $S^{-1}R$ and prime ideals of $R$ disjoint from $S$, $\idealp^{e}$ is of height 1.
$S^{-1}R$ is a Noetherian UFD, so $\idealp^{e} = \frac{a}{s} S^{-1}R = a S^{-1}R$ for some $a \in R$. 
The identity $\idealp = (\idealp^{e})^{c}$ says $aR \subset \idealp$ 
Furthermore, because $a / 1$ is a prime element in $S^{-1}R$, $a$ is a prime element in $R$:
suppose $bc \in aR$. 
Then $\frac{b}{1} \frac{c}{1} \in aS^{-1}R$, and because $a / 1$ is prime, either $b / 1 \in aS^{-1}R$ or $c / 1 \in aS^{-1}R$. 
\sout{It follows that $b \in aR$ or $c \in aR$, proving $a$ is prime in $R$}  \textbf{NOT NECESSARILY TRUE!!!}
Hence $\idealp$ contains a prime ideal $\<a\>$, and the height of $\idealp$ forces $\idealp = \<a\>$. 

In all, any prime ideal in $R$ of height 1 is principal. 
This proves $R$ is a UFD. 

\setcounter{subsection}{19} 
\subsection{} 
Let $f = a_d x^d + \ldots a_0$. 
First, assume $f$ is a unit with $fg = 1$, $g = b_t x^t + \ldots + b_0$. 
Writing $fg = \sum_{k = 0}^{d + t} c_k x^k$, recall 
\begin{equation} \label{eq:4_20_prodcoef} 
  c_k = \sum_{i + j = k} a_i b_j
\end{equation}
\eqref{eq:4_20_prodcoef} indicates $a_0 b_0 = 1$, so $a_0$ is a unit (and so is $b_0$), as desired. 
The following shows that $a_d$ is nilpotent. 
First, by the following induction, 
$$a_{d}^{k + 1} b_{t - k} = 0$$
for all $k \geq 0$. 
It is trivially true for $k = 0$ since $0 = c_{d + t} = a_d b_t$. 
Suppose the result is true up to and including $k - 1$. 
By \eqref{eq:4_20_prodcoef}, 
$$0 = c_{d + t - k} = a_d b_{t - k} + a_{d - 1}b_{t - k + 1} + \ldots + a_{d - k} b_{t}$$
Consider multiplying $a_{d}^{k}$ to the equation. 
The inductive hypothesis says that multiplying by $a_{d}^{k}$ annihilates all the summands on the right hande side except for the first summand $a_{d} b_{t - k}$, and of course the left hand side remains 0. 
Hence $a_{d} b_{t - k} a_{d}^{k} = a_{d}^{k + 1} b_{t - k} = 0$, completing the induction. 
Namely, this result shows 
$$a_{d}^{t + 1} b_0 = 0.$$
Since $b_0$ is a unit, it is not a zerodivisor, which forces $a_{d}^{t + 1} = 0$. 
To see that $a_1, \ldots, a_{d - 1}$ are nilpotent as well, notice the following transformation. 
$a_d$ is nilpotent, so $1 - a_{d}^{t}$ is a unit; 
because $f$ is a unit as well, 
\begin{align*}
  f(1 - a_d^{t}) & = f - a_{d}^{t}(a_{d - 1}x^{d - 1} + \ldots + a_0) \\ 
					 & = a_{d} x^{d} + (1 - a_d^{t})(a_{d - 1}x^{d - 1} + \ldots + a_0)
\end{align*}
is a unit. 
$\text{(unit)} - \text{(nilpotent)}$ is once again a unit, so one concludes from the equality of the first and last expressions above that
$$a_{d - 1}x^{d - 1} + \ldots + a_0$$ 
is a unit. 
Now, inductively applying the same technique as above yields $a_1, \ldots, a_{d - 1}$ as nilpotents. 
This finishes the backward direction. 

The forward direction is much easier. 
If $f = a_d x^d + \ldots + a_1 x + a_0$ with $a_0$ a unit in $R$ and $a_1, \ldots, a_d$ nilpotents in $R$, then $f$ is a sum of a unit with a nilpotent (recall that the nilradical is an ideal). 
The sum of a unit and a nilpotent is a unit, so $f$ is a unit. 

\section{Irreducibility of Polynomials} 
\setcounter{subsection}{2} 
\subsection{}
Let $f(x)$ be a polynomial with even-power terms only and suppose $h(x) \mid f(x)$. 
The map 
$$\Phi: R[x] \rightarrow R[x]; \quad \fnrest{\Phi}{R} = \id, \, x \mapsto -x$$ 
is an automorphism $R[x] \overset{\sim}{\rightarrow} R[x]$. 
Since $f$ only has even-power terms, $\Phi(f(x)) = f(x)$. 
Hence, writing $f(x) = h(x)q(x)$, 
$$f(x) = h(-x) q(-x),$$
which shows $h(-x)$ is a factor of $f(x)$ as well. 

\setcounter{subsection}{6} 
\subsection{} 
Let $g$ be another degree $d$ polynomial whose values agree as that of $g$ at $d + 1$ distinct points of $R$. 
Let $h \defeq f - g$. 
Then $h$ has at most degree $d$ but is 0 at at least $d + 1$ distint points of $R$; 
because $R$ is an integral domain, $h$ must be the zero polynomial. 
In other words, $f = g$; 
$f$ is uniquely determined by its value of any $d + 1$ distinct points of $R$. 

\subsection{} 
As given in the hint, first assume $b_1 = b_2 = \cdots = b_d = 0$ and only $b_0 \neq 0$. 
Define 
$$g_0(x) \defeq (x - a_1)(x - a_2) \cdots (x - a_d)$$ 
and let $c_0 \defeq g_0(a_0)$. 
$g_0(x)$ is a degree $d$ polynomial whose $d$ roots are, by definition, $a_1, \ldots a_d$, so $c_0 \neq 0$ and is thus invertible. 
So one can define $k_0 \defeq c_0^{-1} b_0$. 
Then 
$$f(x) \defeq k_0 g_0(x)$$ 
satisfies $f(a_0) = b_0, \, f(a_1) = f(a_2) = \cdots = f(a_d) = 0$. 

For the general case, proceed as follows: 
for all indices $i \in \{0, \ldots, d\}$ such that $b_i \neq 0$, let 
$$g_i(x) \defeq \frac{1}{x - a_i} \prod_{j = 0}^{d} (x - a_j),$$ 
i.e., $g_i(x) = (x - a_1) \cdots (x - a_{i - 1}) (x - a_{i + 1}) \cdots (x - a_d).$
As in the previous paragraph, one can define 
$$k_i \defeq g_i(a_i)^{-1} b_i.$$
Now, if $i_1, \ldots, i_t \in \{0, \ldots, d\}$ are all the indices such that $b_{i_j} \neq 0$, then 
$$f(x) \defeq k_{i_1} g_{i_1}(x) + k_{i_2}g_{i_2}(x) + \ldots + k_{i_t} g_{i_t}(x)$$
satisfies $f(a_0) = b_0, \, f(a_1) = b_1, \ldots, f(a_d) = b_d$. 
Further, this is the unique polynomial to satisfy the desired property by Q7. 



\end{document}
